\def\topic{Evaluation}
\ifdefined\MyHeaderIncluded
\else
\def\qwertyflag{true}

\ifdefined\MyHeaderIncluded

\else
  \newcommand{\foobar}{FooBar}
  \def\MyHeaderIncluded{true}
% if is till the end of the file

\documentclass[10pt,a4paper]{article}
\usepackage[utf8]{inputenc}
%\usepackage[german]{babel}
%\usepackage[cyrillic]{babel}
\usepackage[T1]{fontenc}
%\usepackage[T2A]{fontenc}
%\usepackage[english, russian]{babel}
\usepackage{amsmath}
\usepackage{amsfonts}
\usepackage{amssymb}
\usepackage{amsthm}
\usepackage{mathtools} %coloneqq

\usepackage{bussproofs} %

%\renewcommand{\thesubsection}{\arabic{subsection}}
%\newtheorem{thm}{Theorem}[section]
%\renewcommand{\thethm}{\arabic{subsection}.\arabic{thm}}

\usepackage{thmtools}
\usepackage{hyperref}
\hypersetup{
    colorlinks=true,
    linkcolor=blue,
    filecolor=magenta,      
    urlcolor=cyan,
    pdftitle={Overleaf Example},
    %pdfpagemode=FullScreen,
    }
%\usepackage{stix-type1}%\turnediota
\theoremstyle{definition}

\newtheorem{constructor}{Constructor}[section]
\newtheorem{iaxiom}{Axiom}
\newtheorem{axiom}{Axiom}
\newtheorem*{axiom*}{Axiom}
\newtheorem{definition}{Definition}[section]
\newtheorem*{definition*}{Definition}
\newtheorem{remark}{Remark}[section]
\newtheorem*{remark*}{Remark}
\newtheorem{recall}{Recall}[section]
\newtheorem*{recall*}{Recall}

\newtheorem{lemma}{Lemma}[section]
\newtheorem{theorem}{Theorem}[section]
\newtheorem*{theorem*}{Theorem}
\newtheorem{note}{Note}[section]
\usepackage{etoolbox}

%\newcommand*{\@Dashv}[2]{%
%  \reflectbox{$\m@th#1#2$}%
%}
\newcommand*{\separation}[3]{% 1)binder, 2)set, 3)condition
  \{#1\in #2: #3\}
}

% https://www.freecodecamp.org/news/simple-regex-tricks-for-beginners-3acb3fa257cb/
%  \\mathcal\{T\} |---> \topology
%  \\mathcal\{([^\}]*)\}   |--->  $1
\newcommand{\AND}{\ \ \&\ \ }
\newcommand{\any}{\mathrm{\ any}}
\newcommand{\borel}{\mathcal{B}}
\newcommand{\can}{\mathrm{can}}
\newcommand{\canonize}{\mathrm{canonize}}
\newcommand{\closure}{\mbox{cl}}
\newcommand{\complex}{\mathbb{C}}
\newcommand{\converges}{\longrightarrow}
\newcommand{\countable}{\mbox{\ countable}}
\newcommand{\ldefined}{\mbox{\ defined}}
\newcommand{\dom}{\mbox{dom}}
\newcommand{\Dom}{\mbox{Dom}}
\newcommand{\Done}{\noindent \textbf{Done.}\\}
\newcommand{\lexists}{\mbox{\ exists}}
\newcommand{\lforall}{\mbox{\ forall}}
\newcommand{\lunique}{\mbox{\ unique}}
\newcommand{\family}{\mathrm{\ family}}
\newcommand{\Fm}{\mathrm{Fm}}
\newcommand{\FuncSymb}{\mathrm{FuncSymb}}
\newcommand{\function}{\mathrm{\ function}}
\newcommand{\functional}{\mbox{\ functional}}
\newcommand{\IFF}{\ \Longleftrightarrow\ }
\renewcommand{\iff}{\ \longleftrightarrow\ }
\newcommand{\IMPLIES}{\ \Longrightarrow\ }
\renewcommand{\implies}{\ \longrightarrow\ }
\newcommand{\Inf}{\mathrm{Inf}}
\newcommand{\injective}{\mbox{\ injective}}
\newcommand{\injection}{\mbox{\ injection}}
\newcommand{\interior}{\mbox{int}}
\newcommand{\intersection}{\bigcap}
\newcommand{\isempty}{\mathrm{\ empty}}
\newcommand{\LB}{\mathrm{LB}}
\newcommand{\nghbset}{\mathcal{N}}
\newcommand{\nonempty}{\mathrm{\ nonempty}}
\newcommand{\Ord}{\mathrm{Ord}}
\newcommand{\power}{\mathcal{P}}
\newcommand{\pr}{\mathrm{pr}}
\newcommand{\fst}{\mathrm{fst}}
\newcommand{\snd}{\mathrm{snd}}
\newcommand{\PredSymb}{\mathrm{PredSymb}}
\newcommand{\Proof}{\noindent \textbf{Proof.}\\}
\newcommand{\ran}{\mathrm{ran}}
\newcommand{\Ran}{\mathrm{Ran}}
\newcommand{\real}{\mathbb{R}}
\newcommand{\relation}{\mathrm{\ relation}}
\newcommand{\set}{\mathrm{\ set}}
\newcommand{\singleton}{\mathrm{\ singleton}}
\newcommand{\semantics}{\mathrm{semantics}}
\newcommand{\semialgebra}{\mathcal{S}}
\newcommand{\Sup}{\mathrm{Sup}}
\newcommand{\surjective}{\mbox{\ surjective\ on\ }}
\newcommand{\Term}{\mathrm{Term}}
\newcommand{\Type}{\mathrm{Type}}
\newcommand{\topology}{\mathcal{T}}
\newcommand{\total}{\mbox{\ total\ on\ }}
\newcommand{\transitive}{\mathrm{\ transitive}}
\newcommand{\union}{\bigcup}
\newcommand{\lTrue}{\mathrm{True}}
\newcommand{\lFalse}{\mathrm{False}}
\newcommand{\lUncertain}{\mathrm{Uncertain}}
\newcommand{\lif}{\mathrm{if}}
\newcommand{\lthen}{\mathrm{then}}
\newcommand{\lelse}{\mathrm{else}}
\newcommand{\lexception}{\mathrm{exception}}
\newcommand{\lundef}{\mathrm{undef}}
\newcommand{\univocal}{\mathrm{\ univocal}}
\newcommand{\univocalat}{\mathrm{\ univocal\ at\ }}
\newcommand{\UB}{\mathrm{UB}}
\newcommand{\Un}{\mathcal{U}}
\newcommand{\Var}{\mathrm{Var}}

\newcommand{\notvarphi}{\phi}
\renewcommand{\phi}{\varphi}

\newcommand{\comment}[1]{\quad\mbox{(#1)}}
% local commands
%\preto\section{\undef\blipo} %etoolbox
\preto\subsection{\undef\blipo} %etoolbox
% https://tex.stackexchange.com/questions/105684/limit-macro-newcommand-scope-to-section
\usepackage{graphicx}
\usepackage{authblk}

\setlength\parindent{0pt} % whole doc of \noindent
\usepackage{soul}
%\usepackage{mathgl}

\fi % guard
\ifdefined\topic
\else
\def\topic{No Titel}
\fi
\title{\topic}
\author[1]{Georgy Dunaev}
\begin{document}
\maketitle
\fi
\section{Evaluation in context of pluralistic semantic}
\begin{definition}[value]
\[
F(x) :\simeq \left(\lif ((F\family)\land(x \in \dom(F)))\ \lthen\ \union\cc{s}{(x,s)\in F}\ \lelse\ \lundef\ \lexception\_\lelse\ \lundef\right)
\]
\end{definition}
\subsection{Value} % of a function.
Definition of the value of the function:
%(in the context of 3--valued logic)
\begin{definition}[semi--equality]
\[
I_{\Fm}(s\approx t) :\IFF 
\]
%(((s \ldefined)\lor(t \ldefined)) \implies (s=t))
\end{definition}

\begin{definition}[equality]?
%\[
%(I^\loc,w_\Fm(s = t)=\ltrue) \IFF \forall \left(v:\prod_{x:Var}w(x)\right),\ I_\Term^{\loc,v}%(t)=I_\Term^{\loc,v}(s)
%\]
%\[
%(I^glob_\Fm(s = t)=\lTrue) \IFF \forall (w:\Var\to\power\Un\setminus\{\emptyclass\} \forall (v:\prod_{x:Var}w(x)),\ I_\Term^{\loc,v}(t)=I_\Term^{\loc,v}(s)
%\]
%(((s \ldefined)\lor(t \ldefined)) \implies (s=t))
\end{definition}

\begin{definition}[quasi--equality](Note: it is good for semantics with partial functions, not the two level set--semantics i use.)
\[
(s\simeq t) :\IFF (((s \ldefined)\lor(t \ldefined)) \implies (s=t))
\]
\end{definition}

\begin{definition}[underdefined class](Note: be careful, since not every undefined class is underdefined, some are overdefined in various ways, so not even semantically equal to eachother.)
\[
\lundef :\simeq (\varepsilon x,\ \bot)
\]
\end{definition}
What if $\varphi$ is assigned with $U$ (i.e. it is assigned with the "uncertain" truth value)?\\
If $(I(\varphi)=U)$, then 
\[
(I((\lif\ \varphi\ \lthen\ t_1\ \lelse\ t_2))=I(\lundef))
\]
should follow.\\
Btw, $I(\lundef)=\emptyclass$.\\
\pagebreak

\subsection{Some simple theorems about classes}
\begin{definition}[some syntactic sugar]
\[
(x\ \mbox{any}) :\iff \top
\]
Context: (x:Var)
%Side condition:\\
%Term $x$ is a variable.
\end{definition}

\begin{definition}[semi--singleton]
\[
(s\ \mbox{ss}) :\iff (!y,\ y\in s)
\]
Side condition:\\
Variable $y$ is not in free variables of term $s$.
\end{definition}

\begin{definition}[semi--singleton subset relation]
\[
(s\ \mbox{sss}\ t) :\iff ((s\ \mbox{ss})\land(s\subseteq t))
\]
\end{definition}

\begin{definition}[function constructor]
\[
(\lambda x,\ t) \ceq \cc{(x,t)}{(x\ \mbox{any})}
\]
Context: ($x$:Var) ($t$:Term)\\
Comment: Term $t$ can contain $x$ as a free variable.\\
Side conditions:\\ 
(variable $x$ is bound by class collection in r.h.s)\\
(variable $x$ is bound by lambda abstraction in l.h.s)\\
%No side conditions.
\end{definition}

\begin{definition}[family constructor]
\[
(\Lambda x,\ t) \ceq \cc{(x,s)}{(s\ \mbox{sss}\ t), x\ \mbox{any}}
\]
% s is a bound variable
Term $t$ can contain $x$ as a free variable.\\
Side conditions:\\
Variables x and s are different.\\ %it should be a normal convention for any definition.
Variable $s$ is not in free variables of $t$.\\
As term $s$ can contain $x$ as a free variable.\\
\end{definition}

\begin{definition*}[LOCAL definition of a value of a function, value of a family, for the next theorem]
\[
f(x) \ceq \bigcup \cc{s}{(x,s)\in f}
\]
\end{definition*}
% \coloneqq
% \vcentcolon = 

\begin{theorem}[evaluation of a function]
Context:\\
(x:Var)(t:Term)(z:Term)
\[
(z\set) \land (t[z/x]\set) \implies
\]
\[
(\lambda x,\ t)(z) = t[z/x]
\]
\end{theorem}
\Proof
To prove the equality, we use set extensionality and MP.
\[
((\lambda x,\ t)(z) \subseteq t[z/x] )
\land
((\lambda x,\ t)(z) \supseteq t[z/x] )
\implies
((\lambda x,\ t)(z) = t[z/x] )
\]
($\subseteq$)\\
Assume 
\[
q\in (\lambda x,\ t)(z)
\]
By definition of $-(-)$ (evaluation),
\[
q\in\union \cc{b}{(z,b)\in(\lambda x,\ t)}.
\]
By destruction of the union,
\[
(\exists s\in \cc{b}{(z,b)\in(\lambda x,\ t)},\ (q\in s)).
\]
Note for the following step: since variable b has no free occurences in t, $t[s/b]=t$.
\[
s\in \cc{b}{(z,b)\in(\lambda x,\ t)}
\IMPLIES
(s\set) \& ((z,s)\in(\lambda x,\ t))
\]
By definition of $\lambda$,
\[
(z,s)\in(\lambda x,\ t)
\IMPLIES
(z,s)\in \cc{(x,t)}{\top}
\]
Unfolding the notation:
\[
(z,s)\in \cc{(x,t)}{\top}
\IMPLIES
(z,s)\in \cc{w}{\exists x,\ w = (x,t) \land \top}
\]

\[
(z,s)\in \cc{w}{\exists x,\ w = (x,t) \land \top}
\IMPLIES
\exists x,\ (z,s) = (x,t) \land \top
\]
Notes: w is free in t, w is different from x. Since otherwise we would write $(x[(z,s)/w],t[(z,s)/w])$.

Take such x. By defining property of ordered pairs,
\[
(z,s) = (x,t)
\IMPLIES
(z=x) \& (s=t)
\]
\[
(q\in s)\& (s=t)
\IMPLIES
q\in t
\]
\[
z=x
\IMPLIES
t[z/x]=t
\]
Therefore
\[
(q\in t) \& (t[z/x]=t)
\IMPLIES
(q \in t[z/x])
\]
\ \\
\ \\
($\supseteq$)\\
Assume $(q \in t[z/x])$.\\
Let $s\ceq t[z/x]$. Then $(q\in s)$.\\
Goal1:$(z,s)\in(\lambda x,\ t)$.\\
\[
x[z/x]=z
\]
\[
(z,t[z/x]) = (x[z/x],t[z/x])=(x,t)[z/x]
\]
Goal1: $(x,t)[z/x]\in\cc{(x,t)}{(x\mbox{ any})}$\\
Goal2: $\exists x,\ (x,t)=((x,t)[z/x])$\\
$x\ceq z$\\
There are no free occurencies of $x$ in $((x,t)[z/x])$, because of the substitution.\\
$c\ceq ((x,t)[z/x])$\\
\[
(((x,t)[z/x])=c)\implies\exists x,\ (x,t)=c
\]
Goal' is done.\\

From the premise, $(x,t)[z/x]$ is a set.\\
Therefore Goal1 is proven.\\
So $(x,t)[z/x]\in (\lambda x,\ t)$.\\
So $(z,t[z/x])\in (\lambda x,\ t)$.\\
So $t[z/x]\in \cc{b}{(z,b)\in(\lambda x,\ t)}$.\\
Since $(q \in t[z/x])$, we also have
\[
q\in \bigcup \cc{b}{(z,b)\in(\lambda x,\ t)}
\]
So $q\in (\lambda x,\ t)(z)$.\\
And we are done.\\
\Done


\begin{theorem}[evaluation of family]
\[
(\Lambda x,\ t)(z) = t[z/x],
\]
and in particular:
\[
(\Lambda x,\ t)(x) = t.
\]
%Context:
%(q term)
\end{theorem}
\Proof
By definitions,
\[
(\Lambda x,\ t)(z) =.
\]


\Done

\subsection{Lambda operatorOLD}
\begin{definition}[lambda operator]
(not good anymore: doesn't work for classes!)
\[
(\lambda x,\ t) = \cc{p}{\exists x,\ p = \langle x,t \rangle}
\labeltagd{def}{lambda operator}
\]
\end{definition}

\ifdefined\qwertyflag
\undef\qwertyflag
\end{document}
\fi